\section*{अध्याय \ref{ch:b2:plant-plasticity} --- पादप अनुकूलन और लक्षणप्ररूपी सुनम्यता}

\begin{solution}{exo:b2:plant-plasticity:1}
सुनम्यता: एक \omterm{def:b2:meiosis-heredity:vocabulary}{जीन-प्ररूप}, और पर्यावरण के अनुसार कई
\omterm{def:b2:meiosis-heredity:vocabulary}{लक्षण-प्ररूप} (एक ही बलूत पर धूप- तथा छाँव-पत्तियाँ)। \omterm{def:b2:plant-plasticity:plasticity}{अनुक्रिया-मानक}:
किसी पर्यावरणीय चर के सापेक्ष किसी \omterm{def:b2:meiosis-heredity:vocabulary}{जीन-प्ररूप} के \omterm{def:b2:meiosis-heredity:vocabulary}{लक्षण-प्ररूप} का
वक्र (प्रकाश के सापेक्ष पत्ती की मोटाई)। \omterm{def:b2:plant-plasticity:plasticity}{अभ्यस्तीकरण}: कोई प्रतिवर्ती
शरीरक्रियात्मक समायोजन (शरद में राई का पाला-कड़ापन)।
\omterm{def:b2:plant-plasticity:plasticity}{अनुकूलन}: पीढ़ियों में वरित कोई वंशागत मेल (किसी कैक्टस का CAM
उपापचय)।
\end{solution}

\begin{solution}{exo:b2:plant-plasticity:2}
एक ओर से प्रकाशित कोई अक्षत प्रांकुर-चोल: प्रकाश की ओर झुकता है। सिरा
किसी अपारदर्शी टोपी से ढका: कोई झुकाव नहीं --- सिरा ही उस
प्रकाश को भाँपता है। सिरा किसी पारदर्शी टोपी से ढका: झुकता है --- वह टोपी
स्वयं हानिरहित है। आधार किसी अपारदर्शी नली से ढका, सिरा खुला: झुकता है ---
संवेदन सिरे पर है और अनुक्रिया उसके नीचे।
\end{solution}

\begin{solution}{exo:b2:plant-plasticity:3}
\omterm{prop:b2:plant-plasticity:sunshade}{धूप-पत्ती}: छोटी, मोटी, दो या तीन खंभ-परतों वाली, मोटे
उपचर्म तथा बहुत रंध्रों वाली, ऊँचे अधिकतम प्रकाशसंश्लेषण और ऊँचे
प्रतिपूरण बिंदु वाली। \omterm{prop:b2:plant-plasticity:sunshade}{छाँव-पत्ती}: बड़ी, पतली, एक खंभ-परत वाली,
प्रति ग्राम अधिक पर्णहरित, नीचा श्वसन, नीचा प्रतिपूरण बिंदु।
वह चुनाव तब होता है जब वह पत्ती कलिका में प्राथमिकांग है, और उस
प्रकाश से जो उसे मिले; कोई परिपक्व पत्ती बदल नहीं सकती।
\end{solution}

\begin{solution}{exo:b2:plant-plasticity:4}
\omterm{prop:b2:plant-plasticity:stomata}{द्वार कोशिकाएँ} पोटैशियम तथा जल लेती हैं, फूलती हैं और, अपनी
भीतरी भित्तियाँ मोटी होने के कारण, धनुष की तरह अलग होकर वह छिद्र खोल देती
हैं। नीला प्रकाश और नीची भीतरी $\mathrm{CO_2}$ उन्हें सुबह खोलते हैं;
\omterm{def:b2:plant-flowering:dormancy}{ऐब्सिसिक अम्ल}, जो सूखती मिट्टी की जड़ों से या उस पत्ती से आता है, उन्हें
किसी सूखे में बंद कर देता है, और बहुत शुष्क वायु भी।
\end{solution}

\begin{solution}{exo:b2:plant-plasticity:5}
$\varphi = R/(R + 0.77\,FR)$, जिसमें $FR = 1$: $1.15$: $0.60$; $0.7$:
$0.48$; $0.2$: $0.21$; $0.05$: $0.06$। देहली $0.4$ उस खुली जगह
($0.48$) और उस छत्र ($0.21$) के बीच पड़ती है: \omterm{thm:b2:plant-plasticity:phytochrome}{छाया-परिहार}
छत्र के नीचे चालू होता है।
\end{solution}

\begin{solution}{exo:b2:plant-plasticity:6}
$E = 0.25\times 0.025 = \qty{6.25}{mmol.m^{-2}.s^{-1}}$; $A =
0.25\times 100/1.6 = \qty{15.6}{\micro mol.m^{-2}.s^{-1}}$; $A/E =
2.5\times 10^{-3}$ मोल कार्बन प्रति मोल जल, यानी प्रति कार्बन 400 जल
के अणु, यानी प्रति ग्राम कार्बन $400\times 18/12 = \qty{600}{g}$
जल। 12 घंटों में: $6.25\times 10^{-3}\times 43200 =
\qty{270}{mol/m^2}$, यानी प्रति वर्ग मीटर \qty{4.9}{L}।
\end{solution}

\begin{solution}{exo:b2:plant-plasticity:7}
$E = 0.25\times 0.005 = \qty{1.25}{mmol.m^{-2}.s^{-1}}$, $A$
अपरिवर्तित: $A/E = 0.0125$, प्रति कार्बन 80 जल, \qty{120}{g/g} ---
यानी दिन की लागत का पाँचवाँ भाग।
\end{solution}

\begin{solution}{exo:b2:plant-plasticity:8}
$v = 2r^{2}\Delta\rho g/9\eta = 2\times 4\times 10^{-12}\times
400\times 9.81/0.09 = \qty{3.5e-7}{m/s}$: यानी \qty{0.35}{\micro m/s},
और उस कोशिका को पार करने में लगभग \qty{30}{s}। किसी क्लाइनोस्टेट पर उस
कोशिका के सापेक्ष गुरुत्व की दिशा बैठने के हर कुछ सेकंडों में बदल जाती
है, इसलिए वे कण कभी एक ओर जमा नहीं होते और उस पौधे के समाकलन-काल (मिनटों)
पर औसत उद्दीपन शून्य है।
\end{solution}

\begin{solution}{exo:b2:plant-plasticity:9}
दो घंटों में माध्य विस्तार $0.04$; दोनों भाग $1.2\times 0.04 = 0.048$
और $0.8\times 0.04 = 0.032$; $\theta = 20\times 0.016/2 = 0.16$
रेडियन, यानी लगभग \qty{9}{\degree}। झुकाव प्रति घंटे $0.08$ रेडियन से
जमा होता है; $\pi/2$ में लगभग \qty{20}{h} लगते हैं।
\end{solution}

\begin{solution}{exo:b2:plant-plasticity:10}
भंडार \qty{50}{cm} तने की छूट देते हैं। बिच्छूबूटी-क्यारी: \qty{30}{mg}
में 7.5 दिनों में \qty{30}{cm} --- वह अंकुर भंडार बचाकर ही
प्रकाश तक पहुँच जाता है। वन: \qty{10}{m} पहुँच से बाहर है; वह अंकुर
अपने \qty{50}{mg} बारह दिनों में आधे मीटर पीले तने पर ख़र्च करके
प्रकाशहीन मर जाता है, जबकि कोई छाया-सहिष्णु रणनीति --- कुछ पतली
पत्तियाँ, धीमी वृद्धि --- उसे वर्षों जीवित रख सकती थी।
\end{solution}

\begin{solution}{exo:b2:plant-plasticity:11}
हिम पहले अंतरकोशिकीय स्थानों में बनता है, जहाँ वह विलयन
अधिक तनु है; हिम के ऊपर वाष्प-दाब उस कोशिका के जल के ऊपर से नीचा है,
इसलिए जल उस कोशिका को छोड़कर हिम से जा मिलता है और वह कोशिका
निर्जलित होकर अपनी झिल्लियाँ ढहा बैठती है। सप्ताह भर की ठंड उस
कोशिका में ऐसी शर्कराएँ तथा रक्षक प्रोटीन भर देती है जो उसका हिमांक
गिराते और झिल्लियाँ स्थिर करते हैं, और ऐसे हिमरोधी प्रोटीन बनाती है जो
रवों की वृद्धि सीमित करते हैं। उस कार्यक्रम को साल भर चालू रखना
कार्बन को वृद्धि से हटाकर रक्षा में लगाता है और विभाजन धीमा करता है:
सदा कड़ा रहने वाला पौधा छोटा पौधा है।
\end{solution}

\begin{solution}{exo:b2:plant-plasticity:12}
दूर-लाल संकेत तब ईमानदार है जब उस पड़ोसी को लाँघा जा सके, और तब झूठ जब
वह कोई वृक्ष हो; वह गर्म सप्ताह अप्रैल में ईमानदार है और फ़रवरी में झूठ;
किसी सघन फ़सल में हर पौधे के पड़ोसी उसके अपने बराबर हैं, इसलिए वह संकेत
रूप में ईमानदार पर व्यर्थ है, और प्रजनक ऐसे पौधे चुनते हैं जो उसकी
उपेक्षा करें। कोई \omterm{def:b2:plant-plasticity:plasticity}{अनुक्रिया-मानक} इस बात का लेखा है कि उस
वंशक्रम के अतीत में हर संकेत ने कितनी बार सच बोला, और वह
जब भी वर्तमान उस अतीत से भिन्न हो तब चूक जाता है।
\end{solution}

\begin{solution}{pb:b2:plant-plasticity:1}
\textbf{1.} खुले में: $1.15/(1.15 + 0.77) = 0.60$; छत्र के नीचे: $0.2/(0.2 +
0.77) = 0.21$।
\textbf{2.} नीचा $\varphi$: ऊपर पौधे हैं --- लंबा होइए। किसी
छाया-कपड़े के नीचे $\varphi$ $0.60$ ही रहता है: यानी कोई अँधेरा दिन, कोई पड़ोसी नहीं; कोई
\omterm{thm:b2:plant-plasticity:phytochrome}{छाया-परिहार} नहीं, केवल प्रकाश के अभाव में धीमी वृद्धि।
\textbf{3.} $1 + 2(0.60 - 0.21) = 1.78$: यानी प्रतिदिन \qty{2.7}{cm}।
\textbf{4.} \qty{21}{cm} के मुक़ाबले \qty{37}{cm}।
\textbf{5.} उस खुली जगह की ओर झुकना (प्रकाशानुवर्तन, नील-प्रकाश
ग्राही फ़ोटोट्रोपिन से) और ऊँचे $R{:}FR$ वाले भाग पर तेज़ दीर्घीकरण
तथा पत्ती का पुनर्विन्यास (फ़ाइटोक्रोम से)।
\textbf{6.} छाँव-पत्तियाँ दुर्लभ प्रकाश सस्ते में बटोरने के लिए बड़ी और
पतली गढ़ी जाती हैं; पूर्ण धूप में वे गरमा जाती हैं, प्रकाश-संदमित होती
और झुलसती हैं, और वह पौधा उन्हें धूप-पत्तियों से बदलता है।
\textbf{7.} \qty{15}{mm} क्षेत्र पर प्रतिदिन \qty{2.7}{cm} का अर्थ है
घंटे भर में \qty{0.11}{cm}, यानी प्रति घंटा $0.074$ का सापेक्ष विस्तार;
छायादार भाग $1.3\times 0.074 = 0.096$, प्रकाशित भाग $0.7\times 0.074 =
0.052$।
\textbf{8.} $\theta = 15\times(0.096 - 0.052)/2 = 0.33$ रेडियन, यानी
घंटे भर में लगभग \qty{19}{\degree}।
\textbf{9.} \qty{60}{\degree} यानी $1.05$ रेडियन: लगभग तीन घंटे।
\textbf{10.} $v = 2\times 6.25\times 10^{-12}\times 400\times
9.81/0.09 = \qty{5.5e-7}{m/s}$, यानी \qty{0.55}{\micro m/s}: उस कोशिका
को पार करने में \qty{22}{s}।
\textbf{11.} प्रकाशानुवर्तन, और \omterm{thm:b2:plant-plasticity:phytochrome}{छाया-परिहार} स्वयं भी उस
कोण को चढ़ा देता है जिस पर वह प्ररोह उगने को राज़ी है: प्रकाश ही वह
संसाधन है जिसके लिए लड़ाई है, और जो तना गुरुत्व के विरुद्ध सीधा हो जाता
वह फिर छाया में लौट जाता।
\textbf{12.} वह जड़ नीचे की ओर तब तक झुकती है जब तक ऊर्ध्वाधर न हो जाए:
स्टैटोलिथ निचले भाग में बैठते हैं, ऑक्सिन वहीं जमा होती है, और किसी जड़
में वह अतिरिक्त ऑक्सिन वृद्धि संदमित करती है, इसलिए ऊपरी भाग तेज़ बढ़ता
है। वही पुनर्वितरण, कोशिकाओं की उलटी संवेदिता।
\textbf{13.} $E = 0.2\times 0.015 = \qty{3}{mmol.m^{-2}.s^{-1}}$;
$2\times 10^{-3}$ \unit{m^2} तथा \qty{50400}{s} पर: \qty{0.30}{mol},
यानी प्रतिदिन \qty{5.4}{g} जल।
\textbf{14.} $A = 0.2\times 80/1.6 = \qty{10}{\micro mol.m^{-2}.s^{-1}}$;
प्रतिदिन प्रति वर्ग मीटर $0.50$ मोल कार्बन; $A/E = 3.3\times
10^{-3}$: प्रति कार्बन 300 जल, \qty{450}{g/g}।
\textbf{15.} $0.50\times 2\times 10^{-3} = \qty{1}{mmol}$: यानी
प्रतिदिन \qty{12}{mg} कार्बन।
\textbf{16.} खुली जगह: $E = \qty{6}{mmol.m^{-2}.s^{-1}}$, प्रतिदिन
\qty{10.9}{g}; दक्षता \qty{900}{g/g}।
\textbf{17.} $E = 0.02\times 0.03 = \qty{0.6}{mmol.m^{-2}.s^{-1}}$
(दस गुना नीचे); $A = 0.02\times 200/1.6 = \qty{2.5}{\micro
mol.m^{-2}.s^{-1}}$ (चार गुना नीचे)। जल अधिक गिरता है: ज्यों-ज्यों वह
पत्ती अपनी भीतरी $\mathrm{CO_2}$ खींच लेती है त्यों-त्यों कार्बन का
प्रवणता-अंतर चौड़ा होता है जबकि जल का नहीं हो सकता।
\textbf{18.} \qty{4.25}{g} जल, जिसमें से \qty{0.85}{g} खोना है;
प्रतिदिन \qty{10.9}{g} ($0.78$ ग्राम प्रति घंटा) पर वह घंटे भर में मुरझा जाता है।
\textbf{19.} CAM धीमा है और उसे मांसल भंडारण ऊतक चाहिए; किसी छत्र के
नीचे प्रकाश के लिए दौड़ता कोई अंकुर, जहाँ जल-हानि वैसे भी कम है,
मितव्ययिता नहीं बल्कि गति चाहता है।
\textbf{20.} $100/2 = \qty{50}{cm}$।
\textbf{21.} बिच्छूबूटी-क्यारी: \qty{80}{mg} में 15 दिनों में \qty{40}{cm} ---
वह प्रकाश तक पहुँच जाता है। वन: \qty{8}{m} अप्राप्य है; वह अपने
भंडार आधे मीटर पर ख़र्च करके मर जाता है। वन में बेहतर: छोटा रहिए,
पतली छाँव-पत्तियाँ बनाइए और किसी खुली जगह की प्रतीक्षा कीजिए --- यानी बीच की राह।
\textbf{22.} भोजपत्र: कोई तीखा मानक, जो नीचे $\varphi$ पर ज़ोर से लंबा
होता है, क्योंकि किसी अग्रगामी के पड़ोसी उसी के आकार के हैं। बीच: कोई
सपाट मानक, थोड़ा दीर्घीकरण और बहुत सहिष्णुता, क्योंकि उसके पड़ोसी
वृक्ष हैं।
\textbf{23.} \omterm{thm:b2:plant-plasticity:phytochrome}{छाया-परिहार} अनुक्रिया निष्क्रिय कीजिए --- फ़ाइटोक्रोम B
को सक्रिय रखिए या जिन \omterm{prop:b2:cell-differentiation:transcriptional}{अनुलेखन कारकों} को वह रोके हुए है उन्हें हटाइए ---
और पुष्पन-काल, अंकुरण तथा तने की मज़बूती जाँचिए, जिन्हें वही
पथ नियंत्रित करता है।
\textbf{24.} जो पौधा पड़ोसी भाँप ही न सके वह हर उस बार लाँघ लिया जाता है
जब वह जीत सकता था; और जो पौधा उस संकेत पर सदा विश्वास कर ले वह हर वृक्ष
के नीचे स्वयं को गँवा देता है। जीतने वाला मानक उस संकेत पर उसी अनुपात
में दाँव लगाता है जिस अनुपात में वह उस वंशक्रम के अतीत में सही
निकला।
\textbf{25.} खुले में $\varphi = 0.60$ और छत्र के नीचे $0.21$;
14 दिनों बाद \qty{37}{cm}; छाया में प्रतिदिन \qty{5.4}{g} जल,
और खुली जगह में \qty{10.9}{g}।
\end{solution}
